\documentclass[10pt,a4paper]{amsart}
\usepackage[margin=19mm]{geometry}
\usepackage{amsmath,amssymb,mathtools}
\usepackage[scr=rsfs,cal=euler]{mathalfa}
\usepackage{xfrac}
\usepackage[unicode,hidelinks,bookmarks=false]{hyperref}
\newcommand{\ie}{i.e.\ }
\newcommand{\cf}{cf.\ }
\newcommand{\resp}{respectively\ }
\newcommand{\Subsection}{\subsection}
\providecommand{\nobreakdash}{\nobreak\hbox{-}}
\setlength{\emergencystretch}{2em}
\multlinegap0pt
\usepackage{mathtools}
\usepackage{amssymb}
\usepackage{xcolor}
\usepackage{cancel}
\usepackage{bm}
\usepackage{tikz}
\newtheorem{proposition}{Proposition}
\newtheorem{lemma}{Lemma}
\newcommand{\G}{\mathcal{G}}
\newcommand\Id{{\mathrm{Id}}}
\newcommand{\Mno}{\mathrm{M}^{d\times d}}
\newcommand\N{\mathbb{N}}
\newcommand{\On}{\mathrm{O}(d)}
\newcommand\R{\mathbb{R}}
\newcommand{\SOn}{\mathrm{SO}(d)}
\newcommand\calU{\mathcal{U}}
\newcommand\dist{{\mathrm{dist}}}
\newcommand\dt{{\,\mathrm{d}t}}
\newcommand{\ol}{\overline}
\title{Phase-field approximation of sharp-interface energies accounting for lattice symmetry}
\author{Sergio Conti, Vito Crismale, Adriana Garroni, Annalisa Malusa}
\date{Source: arXiv:2601.07497v1 (CC BY 4.0)}
\begin{document}
\maketitle

\noindent\emph{Source and licence.} Phase-field approximation of sharp-interface energies accounting for lattice symmetry. Version arXiv:2601.07497v1 (CC BY 4.0), distributed by arXiv under the Creative Commons Attribution 4.0 International licence, http://creativecommons.org/licenses/by/4.0/, as recorded in the arXiv licence metadata for this exact version and shown on the abstract page https://arxiv.org/abs/2601.07497v1. Full licence notice: \texttt{licenses/arxiv-2601.07497-CC-BY-4.0.txt}.\par\smallskip

\noindent\textit{Editorial scope.} Counted unit: the article's proposition prop:proprietag (source0.tex lines 1743--1758). The article's complete original proof of it is delivered with it (source0.tex lines 1759--1911, one \texttt{proof} environment of 153 lines). No mathematics is written, repaired or re-derived: the statement and the proof are the article's own lines, in the article's own notation, and no retained line is substituted or re-flowed.  Retained uncounted as the source-local context the proof needs: lines 351--353; lines 1514--1524; lines 1576--1578; lines 1710--1720.  Nothing was omitted from any retained span, so the package carries no omission record.  No retained line contains a citation, so the wrapper carries no bibliography and no bibliography entry.  No retained line contains a figure environment, an image-inclusion command or drawing code, so no image is delivered and the package declares no asset dependency.  Editorial formatting only, in addition: an AMS \texttt{amsart} preamble that restates the article's own declarations and macros used by the retained lines, namely the environments \texttt{proposition}, \texttt{lemma} on separate counters, together with the standard packages those lines need.  The retained environments print in this document as Proposition 1, Lemma 1 in order of appearance, whereas the article prints the same items with the numbering of its own full document.  The complete native source of the article as distributed in the arXiv e-print for this exact version is bundled unchanged as \texttt{sources/192506/source0.tex}.
\begin{equation}\label{0405230818}
	\ell:=\lim_{s\to 1^-} \frac{1-s}{|\log(1-s)|}f(s)>0, 
\end{equation}

\begin{proposition}\label{p:bpg}
For every $\lambda\in [0,+\infty]$,  $R,\,  R^-,\, R^+ \in \On$, $G\in \G$ it holds that 
\begin{enumerate}
	\item\label{p:bpg:0} $g_\lambda(R, GR)=0$;
	\item\label{p:bpg:G} $g_\lambda( R^-, R^+)=g_\lambda( R^-, GR^+)$;
	\item\label{p:bpg:Id} $g_\lambda( R^-, R^+)=g_\lambda( R^-(R^+)^T, \Id)$;
\item\label{p:bpg:sym} $g_\lambda( R^-, R^+)=g_\lambda(R^+, R^-)$;
\item\label{p:bpg:subadd} $g_\lambda( R^-,R^+)
    \le g_\lambda(R^-, R)+g_\lambda(R,R^+)$, and the same for $g^*$.
\end{enumerate}
\end{proposition}

\begin{proposition}\label{le:equivglambda}
For every $\lambda \in[0,+\infty]$, it holds that $g^*_\lambda=\ol g_\lambda$.
\end{proposition}

\begin{lemma}\label{lemmaprojOn}
 Let $I\subset\R$ be an interval, $\gamma\in W^{1,1}(I;\Mno)$ be such that
 $\dist(\gamma(t),\On)\le 1/2$ almost everywhere.
 Then there exist $C>0$ independent of $\gamma$, and  $\widehat\gamma\in W^{1,1}(I;\On)$
 such that
 \begin{equation}
  |\widehat\gamma'|(t)\le (1+C\dist(\gamma(t),\On)) |\gamma'|(t)\ \text{for a.e.}\ t\in I,\hskip5mm
  \text{ and }\hskip5mm
  \widehat\gamma=\gamma \text{ on } \{\gamma\in \On\}.
 \end{equation}
\end{lemma}

\begin{proposition}\label{prop:proprietag}
For every $\lambda\in [0,+\infty]$, the function $ g^*_\lambda \colon \On \times \On \to [0,+\infty)$ satisfies the following properties:
\begin{enumerate}
\item\label{prop:proprietag:upperbound}  there exists $\kappa \geq 1$ depending only on $d$ and $f$ and such that
	\begin{equation}\label{0105232216}
		g^*_\lambda( R^-, R^+) \leq 1 \wedge  \kappa  \Big(\, | R^+- R^-| \,\big|\log | R^+ - R^-|\big| +| R^+ - R^-|\Big),
	\end{equation}
	for every $ R^+,  R^- \in \On$, $\lambda\in [0,\infty]$;

\item\label{prop:proprietag:RS}  for every $R\in\On$ and for every sequence   $(R_n)_n \subset \On $  converging to  $R$, with $R_n \neq R$, there exists
	\[\lim_{n \to +\infty} \frac{g^*_\lambda(R_n, R)}{|R_n-R| \,\big|\log |R_n-R|\big|}=\frac{\ell}{2};
	\]
\item\label{prop:proprietag:C0}  $g^*_\lambda \in C^0(\On\times \On)$.
\end{enumerate}
All properties are immediately inherited by $g_\lambda$.
\end{proposition}

\begin{proof} 
\noindent \textit{Proof of \ref{prop:proprietag:upperbound}.}
It suffices to prove the bound for $\lambda\in [0,\infty)$, with $\kappa$ not depending on $\lambda$.
For every $ R^-$, $ R^+\in\On$ we set
\[
\ol\beta(t):=
\begin{dcases}
	 R^-, & t\in \left[0,\frac13\right], \\
R^-+(R^+-R^-)(3t-1), &   t\in \left[\frac13,\frac23\right], \\
	 R^+,&  t\in \left[\frac23,1\right];
\end{dcases}
\qquad 
\ol v(t) :=
\begin{dcases}
	1-3t, &  t\in \left[0,\frac13\right], \\
	0, &   t\in  \left[\frac13,\frac23\right], \\
	3t-2, &  t\in  \left[\frac23,1\right] .
\end{dcases}
\]	
Applying Proposition \ref{le:equivglambda}, and by a direct computation, we obtain
\begin{equation}\label{eqglambdale1}
\begin{split}
	g^*_\lambda( R^-, R^+) =& \ol g_\lambda( R^-, R^+) 
	\leq   \int_0^1 \sqrt{(1-\ol v)^2+4\lambda\, \dist^2(\ol\beta,\On)} \sqrt{f^2(\ol v) |\ol\beta'|^2+|\ol v'|^2} \dt\\
	= & \int_{(0,1/3)\cup (2/3,1)}(1-\ol v) |\ol v'| \dt= 1.
\end{split}
\end{equation}
This completes the proof in the case
$|R^+-R^-|\ge 1$.

Assume now
$|R^+-R^-|< 1$.
We shall complete the proof of \eqref{0105232216} using again the fact that $g^*_\lambda=\ol g_\lambda$, and choosing properly the competitors in the optimization problem defining $\ol g_\lambda$.
Using Lemma~\ref{lemmaprojOn} on the affine map joining $R^-$ with $R^+$, whose image is contained in a neighbouhood of $\On$ of size $|R^--R^+|/2<1/2$, we produce a curve 
$\gamma\in C^1([\frac13, \frac23];\On)$ such that
$\gamma(\frac13)=R^-$, $\gamma(\frac23)=R^+$, and
\[
\int_{\frac13}^{\frac23} |\gamma'(t)|\dt\le  |R^--R^+|+C|R^--R^+|^2.
\]  


For $s\in (0,1)$ chosen below, we set
\begin{equation*}
	\beta(t):=\begin{dcases}
		 R^-, \quad & t\in \left[0, \frac13\right],\\
		\gamma(t),\quad & t\in \left[\frac13, \frac23\right],\\
		 R^+, \quad & t\in \left[\frac23,1\right],
	\end{dcases} \qquad
	v(t):=\begin{dcases}
		1-3t s^{1/2}, \quad & t\in \left[0, \frac13\right],\\
		1- s^{1/2}, \quad & t\in \left[\frac13, \frac23\right],\\
		1 + 3 (t-1)  s^{1/2}, \quad & t\in \left[\frac23,1\right].
	\end{dcases}
\end{equation*}
An explicit computation, similar to \eqref{eqglambdale1}, gives
\[
\begin{split}
	\ol g_\lambda( R^-,  R^+)  & \leq \int_{(0,1/3)\cup(2/3,1)}(1-v) |v'| 
	\dt + \int_\frac13^\frac23 s^{1/2} f(1- s^{1/2}) |\gamma'(t)|\, \dt
	\\  
	& =
	s +(|R^--  R^+|+C|R^--  R^+|^2) s^{1/2} f(1-s^{1/2})
\end{split}
\]
and hence choosing $s=|R^--  R^+|$ gives
\begin{equation}\label{1902241215}
	\ol	g_\lambda( R^-,R^+) \leq  | R^- - R^+| +|R^--  R^+|^{3/2}(1+C|R^--  R^+|)  f(1-| R^- - R^+|^{1/2}) .
\end{equation}
The conclusion then follows recalling  
\eqref{0405230818}
and Proposition~\ref{le:equivglambda}.


\smallskip

\noindent \textit{Proof of \ref{prop:proprietag:RS}.} Let $(R_n)_n\subset \On$ converging to $R$ be given. Notice that $R_n$, $R$ belong to the same connected component of $\On$, say $\SOn$, namely $R_n \in \SOn$ for every $n\in\N$, $R\in \SOn$.   We start by proving the lower bound:
\begin{equation}\label{0205230936}
	\liminf_{n\to +\infty}\frac{g_\lambda(R_n, R)}{|R_n-R| \,\big|\log |R_n-R|\big|}\geq \frac{\ell}{2}.
\end{equation}
By monotonicity, it suffices to prove it for $\lambda=0$.
If \eqref{0205230936} does not hold, then there are $\delta>0$ and 
$R_n\in \SOn$, with $s_n:=|R_n-R|\to0$,
and
$(\ol\beta_n,v_n)\in\calU[R_n,R]$ such that
\begin{equation}\label{0205230936b}
\int_0^1 (1-v_n) \sqrt{f^2(v_n) |\ol\beta_n'|^2+|v_n'|^2} \dt \le
 \frac{\ell-\delta}{2} s_n |\log s_n|.
\end{equation}
For every $n$, since
$v_n(0)=v_n(1)=1$ one obtains
\begin{equation}\label{0205231201}
{(1-v_{n}(t))^2} \le
\int_0^1  (1-v_n)|v_n'|\dt \le
 \frac{\ell}{2} s_n |\log s_n|, \qquad t\in (0,1).
\end{equation}
Moreover, using \eqref{0205230936b} and
$\int_0^1  |\ol\beta_n'| \dt \geq s_n$, we get
\begin{equation}\label{eqmin1vnfvn}
 \begin{split}
 \frac{\ell-\delta}{2} s_n |\log s_n|\ge&
\int_0^1 (1-v_n) f(v_n) |\ol\beta_n'|\dt
\ge s_n \min_{[0,1]} (1-v_n) f(v_n).
\end{split}
\end{equation}
Using first \eqref{eqmin1vnfvn} and then
\eqref{0205231201},
\begin{equation}
\begin{split}
 \frac{\ell-\delta}{2} \ge&
 \min_{[0,1]}
\frac{(1-v_n) f(v_n)}{|\log(1-v_n)|}
\frac{|\log(1-v_n)|} {|\log s_n|}
\ge\min_{[0,1]}
\frac{(1-v_n) f(v_n)}
{|\log(1-v_n)|}
\frac{{|\log (\frac\ell2 s_n |\log s_n|)|}}{ 2|\log s_n|}.
\end{split}\end{equation}
Now we take the limit $n\to\infty$.
By \eqref{0205231201},
$v_n\to1$ uniformly. Therefore
(recalling \eqref{0405230818})
the first fraction converges to $\ell$.
At the same time, since $s_n\to0$ the second fraction converges to $1/2$.
This gives a contradiction, since we assumed $\delta>0$.
This proves
\eqref{0205230936} and then the lower bound.


The upper bound follows immediately using first \eqref{1902241215} and then
\eqref{0405230818}. Precisely, with $s_n:=|R_n-R|\to0$,
\[
\begin{split}
	\limsup_{n\to +\infty}  \frac{g_\lambda( R_n, R)}{|R_n-R| \,\big|\log |R_n-R|\big|}
	\leq&
	\limsup_{n\to +\infty}\frac{s_n +
	(1+Cs_n)
	s_n^{3/2} f(1-s_n^{1/2 })}
	{s_n \,\big|\log s_n\big|}
	\\  =&
	\limsup_{n\to +\infty}\frac{s_n^{1/2}
	f(1-s_n^{1/2})}
	{2\big|\log s_n^{1/2}\big|} =
	\frac{\ell}{2}.
\end{split}
\]

Since we used only the bound in
\eqref{1902241215}, which does not depend on $\lambda$, the limit is uniform in $\lambda$, and therefore \ref{prop:proprietag:RS} holds also for $\lambda=\infty$.
\medskip

\noindent \textit{Proof of \ref{prop:proprietag:C0}.} This follows by combining \eqref{0105232216}, which guarantees the continuity at the origin, and the subadditivity stated in
Proposition~\ref{p:bpg}\ref{p:bpg:subadd}.
\end{proof}

\end{document}
